The direct plug-in assertion follows immediately from the defining identity
\[
\tau_P(x_0)=\mu_{1,P}(x_0)-\mu_{0,P}(x_0).
\]
We prove the nontrivial assertion about the oracleized sample estimator.  Put
\[
\eta_P(x)=\pi_P(x)\mu_{1,P}(x)+(1-\pi_P(x))\mu_{0,P}(x),\qquad
\nu_P(x)=\pi_P(x)\{1-\pi_P(x)\}.
\]
For the oracleized normal equations, define
\[
Z_P=\frac{Y-\eta_P(X)}{A-\pi_P(X)},\qquad
W_P=\{A-\pi_P(X)\}^2.
\]
The denominator is bounded away from zero because \(A\in\{0,1\}\) and
\(\varepsilon\leq\pi_P\leq1-\varepsilon\); bounded outcomes give
\(|Z_P|\leq2/\varepsilon\).  Moreover,
\[
\begin{aligned}
\mathbb E_P[W_PZ_P\mid X]
&=\mathbb E_P[(A-\pi_P(X))(Y-\eta_P(X))\mid X]\\
&=\pi_P(X)\{1-\pi_P(X)\}\{\mu_{1,P}(X)-\mu_{0,P}(X)\}\\
&=\nu_P(X)\tau_P(X),\\
\mathbb E_P[W_P\mid X]&=\nu_P(X).
\end{aligned}
\]
Thus the weighted conditional regression function of \(Z_P\) is exactly
\(\tau_P\), not a newly defined estimand.

With \(h_n=n^{-1/(2\gamma+d)}\), the oracleized evaluation equations are
weighted local-polynomial normal equations for this regression, with the
monomial vector through degree \(\lceil\gamma\rceil-1\).  Their population
local Gram matrix is an integral against
\(\nu_P(x)f_P(x)\).  Density-band and overlap bounds give
\[
\varepsilon(1-\varepsilon)f_-\leq
\nu_P(x)f_P(x)\leq f_+/4,
\]
so its eigenvalues are bounded above and below by positive constants depending
only on the fixed polynomial moment matrix and the displayed class constants.
The spectral safeguard in the corrected construction is therefore inactive
whenever the empirical Gram differs from its expectation by less than half
the population lower eigenvalue.  A direct independence and second-moment
calculation, entry by entry in the fixed-dimensional Gram matrix, bounds the
probability of the complementary event by \(C/(nh_n^d)\).

On the good-Gram event, the Taylor polynomial of \(\tau_P\) at \(x_0\) has
uniform remainder at most \(CLh_n^\gamma\).  Substitution in the weighted
normal equations therefore bounds the deterministic approximation part of
the fitted intercept by \(Ch_n^\gamma\).  Conditional on the design and
treatments, the centered summands
\[
W_P\{Z_P-\tau_P(X)\}
\]
are uniformly bounded and have conditional mean zero after averaging over
treatment and outcome at each covariate.  The local Gram bounds then give
variance at most \(C/(nh_n^d)\) for the fitted intercept.  Jensen's inequality
turns this into mean absolute noise at most \(C(nh_n^d)^{-1/2}\).  Truncation
to \([-2,2]\) cannot increase error and makes the bad-Gram contribution at
most \(C/(nh_n^d)\).  Hence
\[
\mathbb E_P|\widehat\tau_n^{\mathrm{or}}(x_0)-\tau_P(x_0)|
\leq C\left\{h_n^\gamma+(nh_n^d)^{-1/2}+(nh_n^d)^{-1}\right\}
\leq Cn^{-\gamma/(2\gamma+d)}.
\]
The calculation uses only the fixed density band, overlap, bounded outcomes,
H\"older CATE condition, and independent sampling, so the constant is uniform
over the stated class and compact smoothness subsets.